\section{Fundamentals}

\subsection{Properties of the Probability Measure}

Let \(\pqty{\Omega, \eventSpace, \Prob}\) be a probability space. We give some useful properties of the probability measure \(\Prob\), other than those on the example sheets.

\begin{proposition}[Countable Subadditivity]
    \label{prop:countable-subadditivity}%
    If \(\pqty{A_n}\) is a collection in \(\eventSpace\), then
    \[
        \Prob\pqty{\bigcup_{n \in \NN} A_n} \leq \sum_{n \in \NN} \Prob\pqty{A_n}.
    \]
\end{proposition}

\begin{proof}
    Define \(B_1 = A_1\), and for \(n \geq 2\), we let
    \[
        B_n = A_n \setminus \pqty{A_1 \cup A_2 \cup \cdots \cup A_{n - 1}}.
    \]

    Then \(\pqty{B_n}\) is a disjoint collection of sets in \(\eventSpace\), and we have
    \[
        \bigcup_{n \in \NN} B_n = \bigcup_{n \in \NN} A_n,
    \]
    so by the countable additivity property (see \autoref{def:probability-measure}, second point), we have
    \[
        \Prob\pqty{\bigcup A_n} = \Prob\pqty{\bigcup B_n} = \sum \Prob\pqty{B_n} \leq \sum \Prob\pqty{A_n}
    \]
    since \(B_n \subseteq A_n\).
\end{proof}

\begin{proposition}[Continuity of Probability Measure]
    \label{prop:continuity-of-probability-measure}%
    For some \emph{increasing} sequence of events \(\pqty{A_n}_{n \in \NN}\) in \(\eventSpace\) (i.e., we have \(A_n \subseteq A_{n + 1}\) for all \(n \in \NN\)), then we have \(\pqty{\Prob\pqty{A_n}}\) is increasing in \(n\), and converges to the limit
    \[
        \Prob\pqty{A_n} \to \Prob\pqty{\bigcup_{n = 1}^{\infty} A_n}.
    \]
\end{proposition}

\begin{proof}
    We define \(B_1 = A_1\), and for \(n \geq 2\), take
    \[
        B_n = A_n \setminus \bigcup_{m = 1}^{n - 1} A_m.
    \]

    Then \(\pqty{B_n}\) are disjoint, and we have
    \[
        \bigcup_{k = 1}^{n} B_k = A_n,
    \]
    so
    \[
        \Prob\pqty{A_n} = \Prob\pqty{\bigcup_{k = 1}^{n} B_k} = \sum_{k = 1}^{n} \Prob\pqty{B_k} \to \sum_{k = 1}^{\infty} \Prob\pqty{B_k}.
    \]

    By the countable additivity property, we have
    \[
        \Prob\pqty{\bigcup_{k = 1}^{\infty} B_k} = \sum_{k = 1}^{\infty} \Prob\pqty{B_k},
    \]
    and we also have
    \[
        \bigcup_{k = 1}^{\infty} B_k = \bigcup_{k = 1}^{\infty} A_k,
    \]
    Therefore
    \[
        \sum_{k = 1}^{\infty} \Prob\pqty{B_k} = \Prob\pqty{\bigcup_{k = 1}^{\infty} A_k},
    \]
    which finishes our proof.
\end{proof}

\begin{corollary}[Countinuity of Probability Measure for Decreasing Events]
    \label{cor:continuity-probability-measure-decreasing-sequence}%
     For some \emph{decreasing} sequence of events \(\pqty{A_n}_{n \in \NN}\) in \(\eventSpace\) (i.e., we have \(A_n \supseteq A_{n + 1}\) for all \(n \in \NN\)), then we have \(\pqty{\Prob\pqty{A_n}}\) is decreasing in \(n\), and converges to the limit
    \[
        \Prob\pqty{A_n} \to \Prob\pqty{\bigcap_{n = 1}^{\infty} A_n}.
    \]
\end{corollary}

\begin{proof}
    We apply \autoref{prop:continuity-of-probability-measure} by taking the complements.
\end{proof}

\begin{remark}
    There is \emph{no} topology involved in play here -- this continuity is different, and the definition we gave above is what is meant for this probability measure.
\end{remark}

We recall that in terms of the cardinality (the cardinality measure on finite sets) that we have the Inclusion--Exclusion Principle. A similar formula holds for two events, \(A, B \in \eventSpace\), which says
\[
    \Prob\pqty{A \cup B} = \Prob\pqty{A} + \Prob\pqty{B} - \Prob\pqty{A \cap B}.
\]

Applying this again to three events \(A, B, C \in \eventSpace\), we have
\begin{align*}
    \Prob\pqty{A \cup B \cup C} &=  \Prob\pqty{\pqty{A \cup B} \cup C}\\
    &= \Prob\pqty{A \cup B} + \Prob\pqty{C} - \Prob\pqty{\pqty{A \cap C} \cup \pqty{B \cap C}} \\
    &= \Prob\pqty{A} + \Prob\pqty{B} + \Prob\pqty{C} - \Prob\pqty{A \cap B} - \Prob\pqty{A \cap C} - \Prob\pqty{B \cap C} + \Prob\pqty{A \cap B \cap C}.
\end{align*}

This leads to \autoref{thm:inclusion-exclusion}, the Inclusion--Exclusion Formula.

\begin{theorem}[Inclusion--Exclusion Formula]
    \label{thm:inclusion-exclusion}%%
    If we have events \(A_1, A_2, \ldots, \eventSpace\), then we have
    \[
        \Prob\pqty{\bigcup_{i = 1}^{n} A_i} = \sum_{k = 1}^{n} \pqty{-1}^{k + 1} \sum_{1 \leq i_1 < \cdots < i_k \leq n} \Prob\pqty{A_{i_1} \cap A_{i_2} \cap \cdots \cap A_{i_k}}.
    \]
\end{theorem}

\begin{proof}
    This holds for \(n = 2\) events. Assume it holds for \(n - 1\) events.

    We have
    \begin{align*}
        \Prob\pqty{\bigcup_{i = 1}^{n} A_i} &= \Prob\pqty{\bigcup_{i = 1}^{n - 1} A_i \cup A_n} \\
        &= \Prob\pqty{\bigcup_{i = 1}^{n - 1} A_i} + \Prob\pqty{A_n} - \Prob\pqty{\bigcup_{i = 1}^{n - 1} A_i \cap A_n} \\
        &= \Prob\pqty{\bigcup_{i = 1}^{n - 1} A_i} + \Prob\pqty{A_n} - \Prob\pqty{\bigcup_{i = 1}^{n - 1} \pqty{A_i \cap A_n}}.
    \end{align*}

    Now, we apply the induction hypothesis to the terms
    \[
        \Prob\pqty{\bigcup_{i = 1}^{n - 1} A_i} \qand \Prob\pqty{\bigcup_{i = 1}^{n - 1} \pqty{A_i \cap A_n}},
    \]
    and hence
    \[
        \Prob\pqty{\bigcup_{i = 1}^{n - 1} A_i} = \sum_{k = 1}^{n - 1} \pqty{-1}^{k + 1} \sum_{1 \leq i_1 < \cdots < i_k \leq n - 1} \Prob\pqty{A_{i_1} \cap A_{i_2} \cap \cdots \cap A_{i_k}},
    \]
    and if we let \(B_i = A_i \cap A_n\),
    \[
        \Prob\pqty{\bigcup_{i = 1}^{n - 1} B_i} = \sum_{k = 1}^{n - 1} \pqty{-1}^{k + 1} \sum_{1 \leq i_1 < \cdots < i_k \leq n - 1} \Prob\pqty{A_{i_1} \cap A_{i_2} \cap \cdots \cap A_{i_k} \cap A_n}.
    \]

    This gives exactly as desired.
\end{proof}

\begin{remark}
    Let \(\Omega\) be a finite set, and \(\pqty{\Omega, \eventSpace, \Prob}\) be a probability space, where \(\Prob\) is the uniform measure
    \[
        \Prob\pqty{A} = \frac{\abs{A}}{\abs{\Omega}}.
    \]

    Taking \(A_1, A_2, \ldots, A_n \in \eventSpace\), applying \autoref{thm:inclusion-exclusion} (Inclusion--Exclusion Formula) to \(\Prob\pqty{\bigcup_{i = 1}^{n} A_i}\), we have
    \[
        \abs{A_1 \cup A_2 \cup \cdots \cup A_n} = \sum_{k = 1}^{n} \pqty{-1}^{k + 1} \sum_{1 \leq i_1 < \cdots < i_n \leq n} \abs{A_{i_1} \cap \cdots \cap A_{i_n}},
    \]
    the Inclusion--Exclusion Principle on cardinalities. So the measure version implies the cardinality one.
\end{remark}

Now, if we truncate the above sums, in the small cases, we see for \(n = 2\)
\[
    \Prob\pqty{A \cup B} \leq \Prob\pqty{A} + \Prob\pqty{B},
\]
and for \(n = 3\),
\[
    \Prob\pqty{A \cup B \cup C} \geq \Prob\pqty{A} + \Prob\pqty{B} + \Prob\pqty{C} - \Prob\pqty{A \cap B} - \Prob\pqty{A \cap C} - \Prob\pqty{B \cap C}
\]
but
\[
    \Prob\pqty{A \cup B \cup C} \leq \Prob\pqty{A} + \Prob\pqty{B} + \Prob\pqty{C}.
\]

This leads to \autoref{prop:bonferroni-inequality}, the Bonferroni Inequalities.

\begin{proposition}[Bonferroni Inequalities]
    \label{prop:bonferroni-inequality}%
    We have
    \[
        \Prob\pqty{\bigcup_{i = 1}^{n} A_i} \begin{cases}
            \leq \sum_{k = 1}^{r} \pqty{-1}^{k + 1} \sum_{1 \leq i_1 < \cdots < i_k \leq n} \Prob\pqty{A_{i_1} \cap A_{i_2} \cap \cdots \cap A_{i_k}}, & r \text{ is odd}, \\
            \geq \sum_{k = 1}^{r} \pqty{-1}^{k + 1} \sum_{1 \leq i_1 < \cdots < i_k \leq n} \Prob\pqty{A_{i_1} \cap A_{i_2} \cap \cdots \cap A_{i_k}}, & r \text{ is even}.
        \end{cases}
    \]
\end{proposition}

\begin{proof}
    The base case is true for \(n = 2\). We assume these inequalities hold for \(n - 1\) events.

    If \(r\) is odd, let \(B_i = A_i \cap A_n\), and we have
    \[
        \Prob\pqty{\bigcup_{i = 1}^{n} A_i} = \Prob\pqty{\bigcup_{i = 1}^{n - 1} A_i} + \Prob\pqty{A_n} - \Prob\pqty{B_1 \cup B_2 \cup \cdots \cup B_{n - 1}}.
    \]

    We apply the induction hypothesis to \(r\) is odd for \(A_i\), giving
    \[
        \Prob\pqty{\bigcup_{i = 1}^{n - 1} A_i} \leq \sum_{k = 1}^{r} \pqty{-1}^{k + 1} \sum_{1 \leq i_1 < \cdots < i_k \leq n - 1} \Prob\pqty{A_{i_1} \cap \cdots \cap A_{i_k}},
    \]
    and to \(r - 1\) is even for \(B_i\), giving
    \begin{align*}
        \Prob\pqty{\bigcup_{i = 1}^{n - 1} B_i} &\geq \sum_{k = 1}^{r - 1} \pqty{-1}^{k + 1} \sum_{1 \leq i_1 < \cdots < i_k \leq n - 1} \Prob\pqty{B_{i_1} \cap \cdots \cap B_{i_k}} \\
        &= \sum_{k = 1}^{r - 1} \pqty{-1}^{k + 1} \sum_{1 \leq i_1 < \cdots < i_k \leq n - 1} \Prob\pqty{A_{i_1} \cap \cdots \cap A_{i_k} \cap A_n}.
    \end{align*}

    Substituting these in finishes the proof for \(r\) is odd, and the proof is similar for \(r\) is even.
\end{proof}

\begin{example}[Counting Surjective Functions]
    \label{ex:counting-surjections}%
    Let
    \[
        \Omega = \Bqty{\functiondomain{f}{\Bqty{1, 2, \ldots, n}}{\Bqty{1, 2, \ldots, m}}},
    \]
    and take the subset
    \[
        A = \set{f \in \Omega}{f \text{ is surjective}}.
    \]

    We want to find \(\abs{A}\). We define
    \[
        A_i = \set{f \in \Omega}{i \notin \img f}.
    \]

    Then we have
    \[
        A = A_1\compl \cap A_2\compl \cap \cdots \cap A_m\compl = \pqty{A_1 \cup A_2 \cup \cdots \cup A_m}\compl.
    \]

    Now, we have
    \[
        \abs{A} = \abs{\Omega} - \abs{A_1 \cup \cdots \cup A_m},
    \]
    and by the Inclusion--Exclusion Principle, we have
    \begin{align*}
        \abs{A_1 \cup \cdots \cup A_m} &= \sum_{k = 1}^{m} \pqty{-1}^{k + 1} \sum_{1 \leq i_1 < \cdots < i_k \leq m} \abs{A_{i_1} \cap A_{i_2} \cap \cdots \cap A_{i_k}} \\
        &= \sum_{k = 1}^{m} \pqty{-1}^{k + 1} \binom{m}{k} \pqty{m - k}^n,
    \end{align*}
    so we have
    \[
        \abs{A} = \sum_{k = 0}^{m} \pqty{-1}^{k} \binom{m}{k} \pqty{m - k}^{n}.
    \]
\end{example}

\begin{example}[Derangements]
    \label{ex:derangements}%
    \emph{Derangements} are permutations with no fixed points. Let the event space be
    \[
        \Omega = \Bqty{\text{permutations of }\Bqty{1, 2, \ldots, n}}
    \]
    and the event of derangements be
    \[
        A = \Bqty{\text{derangements}} = \set{f \in \Omega}{\forall i = 1, 2, \ldots, n: f\pqty{i} \neq i}.
    \]

    What is the probability of \(A\)?

    For every \(i = 1, 2, \ldots, n\), we let
    \[
        A_i = \set{f \in \Omega}{f\pqty{i} = i}.
    \]

    Then, we have
    \[
        A = A_1\compl \cap A_2\compl \cap \cdots \cap A_n\compl = \pqty{\bigcup_{i = 1}^{n} A_i} \compl.
    \]

    Therefore, by the additivity of the probability measure, we have
    \[
        \Prob\pqty{A} = 1 - \Prob\pqty{\bigcup_{i = 1}^{n} A_i}.
    \]

    Using \autoref{thm:inclusion-exclusion} (Inclusion--Exclusion formula), we have
    \[
        \Prob\pqty{\bigcup_{i = 1}^{n} A_i} = \sum_{k = 1}^{n} \pqty{-1}^{k + 1} \sum_{1 \leq i_1 < \cdots < i_k \leq n} \Prob\pqty{A_{i_1} \cap A_{i_2} \cap \cdots \cap A_{i_k}}.
    \]

    We have
    \[
        \Prob\pqty{A_{i_1} \cap A_{i_2} \cap \cdots \cap A_{i_k}} = \frac{\pqty{n - k}!}{n!}
    \]
    since we have \(k\) points fixed in the permutations, and therefore \(\pqty{n - k}\) places free.
    
    Hence,
    \begin{align*}
        \Prob\pqty{\bigcup_{i = 1}^{n} A_i} &= \sum_{k = 1}^{n} \pqty{-1}^{k + 1} \binom{n}{k} \frac{\pqty{n - k}!}{n!} \\
        &= \sum_{k = 1}^{n} \pqty{-1}^{k + 1} \frac{1}{k!}.
    \end{align*}

    Hence, we have
    \[
        \Prob\pqty{A} = 1 - \sum_{k = 1}^{n} \pqty{-1}^{k + 1} \frac{1}{k!} = \sum_{k = 0}^{n} \frac{\pqty{-1}^{k}}{k!}.
    \]

    As \(n \to \infty\), we have the behaviour that
    \[
        \Prob\pqty{A} \to \frac{1}{e}.
    \]
\end{example}

\subsection{Independence and Conditional Probability}

Let \(\pqty{\Omega, \eventSpace, \Prob}\) be a probability space.

\begin{definition}[Independence]
    \label{def:independence}%
    We say two events \(A, B \in \eventSpace\) are \emph{independent}, denoted as \(A \indp B\), if
    \[
        \Prob\pqty{A \cap B} = \Prob\pqty{A} \cdot \Prob\pqty{B}.
    \]

    A \emph{countable collection} of events \(\pqty{A_n}_{n \in \NN}\) is said to be \emph{independent}, if for any finite collection of indices \(I = \Bqty{i_1, i_2, \ldots, i_k} \subseteq \NN\), we have
    \[
        \Prob\pqty{\bigcap_{i \in I} A_i} = \prod_{i \in I} \Prob\pqty{A_i}.
    \]
\end{definition}

\begin{remark}
    Pairwise independence \textbf{does not imply} independence.
\end{remark}

\begin{example}[Pairwise Independence does not imply Independence]
    \label{ex:pairwise-indep-does-not-imply-indep}
    We toss a fair coin twice. The sample space is
    \[
        \Omega = \Bqty{\pqty{0, 0}, \pqty{0, 1}, \pqty{1, 0}, \pqty{1, 1}}
    \]
    equipped with probability measure
    \[
        \Prob\pqty{\omega} = \frac{1}{4}
    \]  
    for all \(\omega \in \Omega\). Take
    \[
        A = \Bqty{\pqty{0, 0}, \pqty{0, 1}}, B = \Bqty{\pqty{0, 0}, \pqty{1, 0}}, C = \Bqty{\pqty{1, 0}, \pqty{0, 1}}.
    \]

    We have
    \[
        \Prob\pqty{A} = \Prob\pqty{B} = \Prob\pqty{C} = \frac{1}{2},
    \]
    and
    \[
        \Prob\pqty{A \cap B} = \Prob\pqty{A \cap C} = \Prob\pqty{B \cap C} = \frac{1}{4},
    \]
    so \(\Bqty{A, B, C}\) are \emph{pairwise} independent. But they are not \emph{independent}, since
    \[
        \Prob\pqty{A \cap B \cap C} = \Prob\pqty{\emptyset} = 0.
    \]
\end{example}

\begin{proposition}[Independence implies Independence of Complement]
    \label{prop:indep-compl}%
    If \(A \indp B\), then \(A \indp B\compl\).
\end{proposition}

\begin{proof}
    Investigating the probability, we have
    \begin{align*}
        \Prob\pqty{A \cap B\compl} &= \Prob\pqty{A} - \Prob\pqty{A \cap B} \\
        &= \Prob\pqty{A} - \Prob\pqty{A} \cdot \Prob\pqty{B} \\
        &= \Prob\pqty{A} \pqty{1 - \Prob\pqty{B}} \\
        &= \Prob\pqty{A} \Prob\pqty{B\compl}.
    \end{align*}
\end{proof}

An idea closely related is conditional probability.

\begin{definition}[Conditional Probability]
    \label{def:conditional-prob}%
    Let \(B \in \eventSpace\) with non-zero measure, \(\Prob\pqty{B} > 0\).

    For any \(A \in \eventSpace\), we define the conditional probability of \(A\) given \(B\), is given by
    \[
        \Prob\pqty{\Conditional{A}{B}} = \frac{\Prob\pqty{A \cap B}}{\Prob\pqty{B}}.
    \]
\end{definition}

If \(A \indp B\), then we have
\[
    \Prob\pqty{\Conditional{A}{B}} = \Prob\pqty{A}.
\]

\begin{proposition}[Countable Additivity Property for Conditional Probability]
    \label{prop:countable-additivity-conditional}%
    Let \(\pqty{A_n}_{n \in \NN}\) be a disjoint sequence of events in \(\eventSpace\). Then we have
    \[
        \Prob\pqty{\Conditional{\bigcup_n A_n}{B}} = \sum_{n} \Prob\pqty{\Conditional{A_n}{B}}.
    \]
\end{proposition}

\begin{proof}
    By definition, we have
    \begin{align*}
        \Prob\pqty{\Conditional{\bigcup_{n} A_n}{B}} &= \frac{\Prob\pqty{\bigcup_n A_n \cap B}}{\Prob\pqty{B}} \\
        &= \frac{\Prob\pqty{\bigcup_n \pqty{A_n \cap B}}}{\Prob\pqty{B}} \\
        &= \sum_{n} \frac{\Prob\pqty{A_n \cap B}}{\Prob\pqty{B}} \\
        &= \sum_{n} \Prob\pqty{\Conditional{A_n}{B}}.
    \end{align*}
\end{proof}

\begin{theorem}[Law of Total Probability]
    \label{thm:law-total-probability}%
    Suppose that the events \(\pqty{B_n}_{n \in \NN}\) is a \emph{partition} of \(\Omega\), i.e. pairwise disjoint with union \(\Omega\), with positive measure. i.e. \(\Prob\pqty{B_n} > 0\).

    Let \(A \in \eventSpace\). We have
    \[
        \Prob\pqty{A} = \sum_{n} \Prob\pqty{\Conditional{A}{B_n}} \Prob\pqty{B_n}.
    \]
\end{theorem}

\begin{proof}
    We have
    \begin{align*}
        \Prob\pqty{A} &= \Prob\pqty{A \cap \bigcup_{n} B_n} \\
        &= \Prob\pqty{\bigcup_n \pqty{B_n \cap A}} \\
        &= \sum_{n} \Prob\pqty{B_n \cap A} \\
        &= \sum_{n} \Prob\pqty{\Conditional{A}{B_n}} \Prob\pqty{B_n}.
    \end{align*}
\end{proof}

\begin{theorem}[Bayes' Formula]
    \label{thm:bayes}%
    Let \(\pqty{B_n}\) be a partition of \(\Omega\) with positive measure. We have
    \[
        \Prob\pqty{\Conditional{B_n}{A}} = \frac{\Prob\pqty{\Conditional{A}{B_n}} \cdot \Prob\pqty{B_n}}{\sum_{k} \Prob\pqty{\Conditional{A}{B_k}} \cdot \Prob\pqty{B_k}}.
    \]
\end{theorem}

\begin{proof}
    Indeed, we have
    \begin{align*}
        \Prob\pqty{\Conditional{B_n}{A}} &= \frac{\Prob\pqty{B_n \cap A}}{\Prob\pqty{A}} \\
        &= \frac{\Prob\pqty{\Conditional{A}{B_n}} \cdot \Prob\pqty{B_n}}{\Prob\pqty{A}},
    \end{align*}
    and we apply \autoref{thm:law-total-probability} (Law of Total Probability) for the denominator.
\end{proof}

This formula is the basis of Bayesian statistics. If we know the probability of \(\Prob\pqty{B_n}\), and we have a model which gives us \(\Prob\pqty{\Conditional{A}{B_k}}\), with this formula we can find the posterior probability of \(B_n\) given that \(A\) occurs.

\begin{example}[False Positives for a Rare Condition]
    \label{ex:false-positives}%
    Suppose that a condition A affects \(0.1\%\) of the population. We have a medical test which is positive for \(98\%\) of the affected population, and \(1\%\) of those unaffected.

    Pick an individual at random. What is the probability they suffer from A given they tested positive?

    We define the events
    \[
        A = \Bqty{\text{individuals suffering from A}},
    \]
    and
    \[
        P = \Bqty{\text{individuals tests positive}}.
    \]

    We want to find \(\Prob\pqty{\Conditional{A}{P}}\). The given conditions state that
    \[
        \Prob\pqty{A} = 0.001, \quad \Prob\pqty{\Conditional{P}{A}} = 0.98, \qand \Prob\pqty{\Conditional{P}{A\compl}} = 0.01.
    \]

    Hence, using \autoref{thm:bayes} (Bayes' Formula), we have
    \begin{align*}
        \Prob\pqty{\Conditional{A}{P}} &= \frac{\Prob\pqty{\Conditional{P}{A}} \cdot \Prob\pqty{A}}{\Prob\pqty{\Conditional{P}{A}} \cdot \Prob\pqty{A} + \Prob\pqty{\Conditional{P}{A\compl}} \cdot \Prob\pqty{A\compl}} = 0.089 \approx 0.09
    \end{align*}
    so
    \[
        \Prob\pqty{\Conditional{A}{P}} \approx 0.09.
    \]

    This is counterintuitive. We have that
    \[
        \Prob\pqty{\Conditional{A}{P}} = \frac{1}{1 + \frac{\Prob\pqty{\Conditional{P}{A\compl}} \cdot \Prob\pqty{A\compl}}{\Prob\pqty{\Conditional{P}{A}} \cdot \Prob\pqty{A}}}.
    \]

    We have \(\Prob\pqty{A\compl}\) and \(\Prob\pqty{\Conditional{P}{A}}\) are both very close to \(1\), so the relevant ratio simplifies to \(\frac{\Prob\pqty{\Conditional{P}{A\compl}}}{\Prob\pqty{A}}\). But
    \[
        \Prob\pqty{\Conditional{P}{A\compl}} \gg \Prob\pqty{A}
    \]
    and this makes \(\Prob\pqty{\Conditional{A}{P}}\) is \emph{very small}.

    Let's put in some concrete numbers. Suppose we have a population of \(1000\) people, with \(1\) suffering from A. Among the \(999\) not suffering, about \(10\) people will test positive. In total, there will be \(11\) people testing positive. Among those we pick one at random. So the probability of being positive, given they actually test positive, is only \(\frac{1}{11}\).
\end{example}

\begin{example}[Children and Thursday]
    \label{ex:children-thursday}%
    \begin{enumerate}
        \item What is the probability that I have \(2\) boys, given I have \(2\) children, one of whom is a boy?
    
            Since no further information is given, we take all outcomes to be equally likely.

            We denote the events
            \[
                \mathrm{BG} = \Bqty{\text{elder child is a boy, younger is a girl}},
            \]
            and the events \(\mathrm{GB}, \mathrm{GG}, \mathrm{BB}\). We want to find
            \[
                \Prob\pqty{\Conditional{\mathrm{BB}}{\mathrm{BB} \cup \mathrm{BG} \cup \mathrm{GB}}} = \frac{\Prob\pqty{\mathrm{BB}}}{\Prob\pqty{\mathrm{BB} \cup \mathrm{BG} \cup \mathrm{GB}}} = \frac{\frac{1}{4}}{\frac{3}{4}} = \frac{1}{3}.
            \]

        \item What is the probability that I have \(2\) boys, given the eldest one is a boy?
        
            We have
            \[
                \Prob\pqty{\Conditional{\mathrm{BB}}{\mathrm{BB}} \cup \mathrm{BG}} = \frac{1}{2}.
            \]

        \item What is the probability that I have \(2\) boys, given that one of them is a boy born on a Thursday?
        
        Let
        \[
            \mathrm{TG} = \Bqty{\text{elder boy is a boy on Thursday, younger is a girl}},
        \]
        and \(\mathrm{GT}\) similarly, and
        \[
            \mathrm{TN} = \Bqty{\text{elder boy is a boy on Thursday, younger is a boy not born on Thursday}},
        \]
        and \(\mathrm{NT, TT}\) similarly. We want to find
        \begin{align*}
            &\phantom{=} \Prob\pqty{\Conditional{\mathrm{TT} \cup \mathrm{TN} \cup \mathrm{NT}}{\mathrm{TT} \cup \mathrm{TG} \cup \mathrm{GT} \cup \mathrm{TN} \cup \mathrm{NT}}} \\
            &= \frac{\Prob\pqty{\mathrm{TT} \cup \mathrm{TN} \cup \mathrm{NT}}}{\Prob\pqty{\mathrm{TT} \cup \mathrm{TG} \cup \mathrm{GT} \cup \mathrm{TN} \cup \mathrm{NT}}} \\
            &= \frac{\frac{1}{2} \cdot \frac{1}{7} \cdot \frac{1}{2} \cdot \frac{1}{7} + \frac{1}{2} \cdot \frac{1}{7} \cdot \frac{1}{2} \cdot \mathbf{\frac{6}{7}} + \frac{1}{2} \cdot \mathbf{\frac{6}{7}} \cdot \frac{1}{2} \cdot \frac{1}{7}}{\frac{1}{2} \cdot \frac{1}{7} \cdot \frac{1}{2} \cdot \frac{1}{7} + \frac{1}{2} \cdot \frac{1}{7} \cdot \frac{1}{2} + \frac{1}{2} \cdot \frac{1}{7} \cdot \frac{1}{2} + \frac{1}{2} \cdot \frac{1}{7} \cdot \frac{1}{2} \cdot \mathbf{\frac{6}{7}} \cdot 2}
 \\
            &= \frac{13}{27}.
        \end{align*}
    \end{enumerate}
\end{example}

\begin{table}[ht!]
    \centering
    \begin{tabular}{cccc}
        All Applicants & Admitted & Rejected & \% Admitted \\
        \hline
        State & 25 & 25 & 50\% \\
        Independent & 28 & 22 & 56\%
    \end{tabular}
    
    \vspace{1em}

    \begin{tabular}{cccc}
        London Schools & Admitted & Rejected & \% Admitted \\
        \hline
        State & 15 & 22 & 41\% \\
        Independent & 5 & 8 & 38\%
    \end{tabular}
    
    \vspace{1em}

    \begin{tabular}{cccc}
        Cambridge Schools & Admitted & Rejected & \% Admitted \\
        \hline
        State & 10 & 3 & 77\% \\
        Independent & 23 & 14 & 62\%
    \end{tabular}

    \caption{Data of Admissions}
    \label{tab:simpson-paradox-data}
\end{table}

\begin{example}[Simpson's Paradox]
    \label{ex:simpson-paradox}%
    We refer to \autoref{tab:simpson-paradox-data}. This phenomenon is called \emph{cofounding} in statistics. It arises when we aggregate data from disparate populations.

    Let
    \begin{align*}
        A &= \Bqty{\text{individual is admitted}}, \\
        B &= \Bqty{\text{individual is from London}}, \\
        B\compl &= \Bqty{\text{individual is from Cambridge}}, \\
        C &= \Bqty{\text{state school}}, \\
        C\compl &= \Bqty{\text{independent school}}.
    \end{align*}

    We notice that from the latter two tables, that
    \begin{align*}
        \Prob\pqty{\Conditional{A}{B \cap C}} &> \Prob\pqty{\Conditional{A}{B \cap C \compl}}, \\
        \Prob\pqty{\Conditional{A}{B\compl \cap C}} &> \Prob\pqty{\Conditional{A}{B\compl \cap C \compl}}.
    \end{align*}

    However, we have from the first table, that
    \[
        \Prob\pqty{\Conditional{A}{C}} < \Prob\pqty{\Conditional{A}{C\compl}}.
    \]

    Observe that from countable additivity of conditional probability,
    \begin{align*}
        \Prob\pqty{\Conditional{A}{C}} &= \Prob\pqty{\Conditional{A \cap B}{C}} + \Prob\pqty{\Conditional{A \cap B\compl}{C}} \\
        &= \frac{\Prob\pqty{A \cap B \cap C}}{\Prob\pqty{C}} + \frac{\Prob\pqty{A \cap B\compl \cap C}}{\Prob\pqty{C}} \\
        &= \frac{\Prob\pqty{\Conditional{A}{B \cap C}} \cdot \Prob\pqty{B \cap C}}{\Prob\pqty{C}} + \frac{\Prob\pqty{\Conditional{A}{B \compl \cap C}} \cdot \Prob\pqty{B\compl \cap C}}{\Prob\pqty{C}} \\
        &= \Prob\pqty{\Conditional{A}{B \cap C}} \cdot \Prob\pqty{\Conditional{B}{C}} + \Prob\pqty{\Conditional{A}{B\compl \cap C}} \cdot \Prob\pqty{\Conditional{B \compl}{C}} \\
        &> \Prob\pqty{\Conditional{A}{B \cap C\compl}} \cdot \Prob\pqty{\Conditional{B}{C}} + \Prob\pqty{\Conditional{A}{B\compl \cap C\compl}} \cdot \Prob \pqty{\Conditional{B\compl}{C}}.
    \end{align*}

    Now, note that if \(\Prob\pqty{\Conditional{B}{C}} = \Prob\pqty{\Conditional{B}{C\compl}}\), which is not valid here, then we would get
    \begin{align*}
        \Prob\pqty{\Conditional{A}{C}} &> \Prob\pqty{\Conditional{A}{B \cap C\compl}} \cdot \Prob\pqty{\Conditional{B}{C}} + \Prob\pqty{\Conditional{A}{B\compl \cap C\compl}} \cdot \Prob \pqty{\Conditional{B\compl}{C}} \\
        &= \Prob\pqty{\Conditional{A}{B \cap C\compl}} \cdot \Prob\pqty{\Conditional{B}{C\compl}} + \Prob\pqty{\Conditional{A}{B\compl \cap C\compl}} \cdot \Prob \pqty{\Conditional{B\compl}{C\compl}} \\
        &= \Prob\pqty{\Conditional{A}{C\compl}}
    \end{align*}
    hence this does not hold.

    Effectively, we are dealing with inequalities
    \[
        \frac{a_1}{a_2} > \frac{b_1}{b_2} \qand \frac{c_1}{c_2} > \frac{d_1}{d_2},
    \]
    but we cannot say anything about
    \[
        \frac{a_1 + c_1}{a_2 + c_2} \qand \frac{b_1 + d_1}{b_2 + d_2}.
    \]
\end{example}